Se reescribió el proceso ETL de forma modular y reproducible. Los arribos pesqueros se ingieren con granularidad de una fila por (fecha, especie, unidad económica) uniendo dos fuentes complementarias: el export histórico de COBI (2016-2021) y los avisos de cosecha (``AVISOS MAYORES/MENORES DE COSECHA'') de CONAPESCA (2022-2026). La unión se resuelve por prioridad temporal ---COBI hasta el 31-diciembre-2021 y CONAPESCA a partir de 2022, cada archivo acotado a su propio año--- de modo que no hay doble conteo en la costura de 2021. El lector de CSV autodetecta el encabezado (el preámbulo de los avisos varía entre 2 y 4 líneas) y la limpieza acepta rangos de fecha explícitos. El conjunto consolidado cubre así 2017-2026, casi duplicando las temporadas disponibles de langosta en San Quintín (de $\sim$5 a $\sim$9). Como fuente de temperatura superficial del mar (SST) se adoptó el producto reprocesado L4 de Copernicus Marine Service (OSTIA, \texttt{METOFFICE-GLO-SST-L4-REP-OBS-SST}), con cobertura diaria global desde 1981; se descargó la serie 1982-2025 y se promedió sobre el bounding box costero de cada unidad económica. Esta fuente reemplaza al cliente \texttt{motuclient} —descontinuado en 2024— por el nuevo SDK \texttt{copernicusmarine}, y provee la profundidad histórica necesaria para construir una climatología de referencia de 30 años.
